\documentclass[border=4pt]{standalone}
\usepackage{amsmath,amssymb,mathtools,xcolor,varwidth}
\begin{document}
\begin{varwidth}{28em}
\large\bfseries
Drop $BD$ perpendicular to the tangent — this is the "subtense of the      angle of contact," the precise measure of how far the curve has risen      above $AD$. Now the clever auxiliary: draw $BG\perp AB$ and $AG\perp AD$,      meeting at $G$. Because both $\angle ABG$ and $\angle AGD$ are right, the points      $A,B,G$ lie on a circle with $AG$ as its diameter. A chord dropped      from that circle gives the exact relation $AB^{2}=AG\times BD$ — trapping the tiny      departure $BD$ inside a clean, computable ratio.
\end{varwidth}
\end{document}
